\documentclass[11pt]{article}
\usepackage[utf8]{inputenc}
\usepackage{amsmath,amssymb}
\usepackage[margin=1in]{geometry}
\usepackage{hyperref}
\emergencystretch=3em

\title{The general power-step tower:\\
$Z_d(c)=\Big(\prod_{i<d}k_i^{-1}\Big)D_d\,c^{-q_{d-1}}\,G_d\big(cb^{S_d}\big)$ for every $(k,h)$}
\author{Claude, with Daniel Murfet}
\date{2026-09-06}

\begin{document}
\maketitle
\sloppy

\begin{abstract}
All exact reductions so far (units~44, 50, 52) are instances of one recursion. With
$q_i=(h_i+1)/k_i$ and $q_{-1}=0$, define $G_0=f$, $G_{j+1}(y)=\int_0^yx^{q_j-q_{j-1}-1}G_j(x)\,dx$ and
$D_0=1$, $D_{d+1}=D_d\,(b^{S_d})^{q_{d-1}-q_d}$ with $S_d=\sum_{i<d}k_i$. Then for every $d$ and
every exponent vector, $Z_d(c)=(\prod_{i<d}k_i^{-1})\,D_d\,c^{-q_{d-1}}\,G_d(cb^{S_d})$. The
all-equal tower is \texttt{iterDivPrim}$(F_{p-1})$ and the second level with exponents $(q,p)$ is
unit~52's noncritical base. The formula was checked by three-dimensional quadrature with three
distinct exponents to $5\cdot10^{-14}$. Zero \texttt{sorry}/\texttt{axiom}.
\end{abstract}

\section{Setting}
Each coordinate contributes a chart power substitution and a scaling; the only data that changes
from level to level is the exponent $q_d$ of the new coordinate versus the exponent $q_{d-1}$ carried
by the previous level, which appear as the power $x^{q_d-q_{d-1}-1}$ in the next integrand and as the
factor $B^{q_{d-1}-q_d}$ in the scale. The three earlier reductions correspond to
$q_d=q_{d-1}$ (power $x^{-1}$: a logarithm), $q_d<q_{d-1}$ (a convergent integral) and the base
level $q_{-1}=0$ (the weighted primitive).

\section{The things formalised}
{\small
\begin{verbatim}
noncomputable def powStep (s : ℝ) (G : ℝ → ℝ) (y : ℝ) : ℝ := ∫ x in Ioc 0 y, x ^ s * G x
noncomputable def genTower (β a : ℝ) (k h : ℕ → ℕ) : ℕ → ℝ → ℝ
  | 0 => quadKernel β a
  | j + 1 => powStep (chartExp k h j - prevExp k h j - 1) (genTower β a k h j)
theorem iterChartGen_step_general ... (hZ : ∀ c, 0 < c →
      iterChartGen β a b k h d c = K * c ^ (-q') * D * G (c * B)) (c : ℝ) (hc : 0 < c) :
    iterChartGen β a b k h (d + 1) c
      = K * (1 / (k d : ℝ)) * c ^ (-(chartExp k h d)) * (D * B ^ (q' - chartExp k h d))
        * powStep (chartExp k h d - q' - 1) G (c * (B * b ^ (k d)))
theorem iterChartGen_eq_genTower ... (d : ℕ) (c : ℝ) (hc : 0 < c) :
    iterChartGen β a b k h d c
      = (∏ i ∈ Finset.range d, (1 : ℝ) / k i) * genScale b k h d * c ^ (-(prevExp k h d))
        * genTower β a k h d (c * b ^ (∑ i ∈ Finset.range d, k i))
theorem genTower_eq_iterDivPrim ... :
    genTower β a k h (d + 1) = iterDivPrim (weightedPrimitive β a (p - 1)) d
theorem genTower_two_eq_noLogBase ... : genTower β a k h 2 = noLogBase β a p q
\end{verbatim}
}

\section{Proof strategy}
The general step is unit~52's step lemma with the exponent of the previous level freed: after the
chart power substitution the integrand is $s^{q-1}s^{-q'}G(cBs)=s^{q-q'-1}G(cBs)$, the scaling lemma
with $\gamma=q-q'-1$ produces $(cB)^{-(q-q')}$, and $c^{-q'}c^{-(q-q')}=c^{-q}$ leaves $B^{q'-q}$ as
the new scale factor. The exact reduction is an induction on $d$ generalising $c$; the base case
$d=0$ is trivial because all conventions ($q_{-1}=0$, $D_0=1$, $S_0=0$) make $Z_0=f$ fit the general
shape. The two identifications are inductions in which the power $x^{-1}$ is rewritten as division
and the definitions unfold by \texttt{rfl}.

\section{Roadblocks and resolutions}
\begin{itemize}
\item The final goal of the induction step differed only by the order of factors; \texttt{ring}.
\item Unfolding \texttt{weightedPrimitive} under a binder with \texttt{rw} fails; the goal was
closed by \texttt{rfl} since all three definitions unfold definitionally.
\end{itemize}

\section{What was Mathlib, what was new}
Mathlib: \texttt{Real.mul\_rpow}, \texttt{Real.rpow\_add}, \texttt{Finset.prod\_range\_succ},
\texttt{Finset.sum\_range\_succ}. New: the general tower, its exact reduction for arbitrary
$(k,h)$, and the identifications with the earlier special towers.

\section{Lessons}
The right generality for an exact reduction is the one in which the base case is trivial: with
$q_{-1}=0$ and $D_0=1$ the kernel itself is a tower of height zero, and no special treatment of the
innermost coordinate is needed.

\section{Outcome}
The chart standard integral with constant $\xi$ is now reduced exactly, for every exponent vector, to
a one-dimensional tower $G_d$. The remaining analytic problem is the growth of $G_d(y)$ as
$y\to\infty$, expected to be $y^{q_{d-1}-q_{\min}}(\log y)^{m-1}$ with $m$ the multiplicity of the
minimal exponent; this would give the full ``only minimal exponents count'' theorem.
\end{document}
